\chapter{Installation de Tockem SPE}

Le logiciel Tockem SPE est une application web développée par le volet EHA de l'association Tockem dans le but de numériser le suivi technique et financier des activités liées aux différents AEP qu'administre le volet. Dans cette partie du document, nous verrons comment installer et utiliser l'application.

Une application web est un logiciel comme on en voit usuellement, à la différence qu'elle s'exécute dans un navigateur. Pour ce qui est du navigateur, il s'agit d'un logiciel qui permet aux individus de consulter des pages sur Internet. Dans notre cas, il ne sera pas nécessaire d'avoir une connexion Internet. Cependant, pour le bon fonctionnement du logiciel, un ensemble d'outils sera nécessaire pour permettre au logiciel de fonctionner sans Internet. Dans notre cas, il est question de \textbf{WAMP}.

\textbf{WAMP} (Windows, Apache, MySQL, PHP) est un ensemble formé d'un serveur web (Apache) qui permet de construire les pages qui vont être affichées dans le navigateur, d'une base de données (MySQL) qui stocke les informations des différentes applications web, et d'un langage de programmation (PHP) qui est utilisé pour créer les différentes pages.

Par la suite, il faut configurer l'application web. C'est à un développeur web de s'en charger~:
\begin{itemize}
    \item Ajout des sources de l'application web dans WAMP
    \item Création et configuration de la base de données de Tockem SPE
    \item Définition d'un fichier de démarrage rapide avec l'extension \texttt{.bat}
\end{itemize}

Il est à noter que ces opérations ne sont effectuées qu'une seule fois, lors de l'installation du logiciel.

Pour exécuter le logiciel, il faut lancer le fichier \texttt{tokem\_spe.bat}, qui va se charger de lancer WAMP et d'ouvrir une fenêtre du navigateur qui affichera l'application.
